%!TEX program = xelatex
\documentclass[12pt, b5paper]{article}
\usepackage{ctex}

\usepackage{geometry}
\geometry{b5paper,left=2cm,right=1cm,top=2.5cm,bottom=2.5cm}
\usepackage{chemformula} %化学公式宏包
\usepackage[version=4]{mhchem} %化学公式宏包
\usepackage{siunitx} %数字，单位宏包
\sisetup{
	separate-uncertainty = true,
	inter-unit-product = \ensuremath{{}\cdot{}}
} %单位间加小圆点
\usepackage{tikz} %Tikz绘图包
\usetikzlibrary{cd}
\usetikzlibrary{chains}
\usetikzlibrary{arrows}
\usetikzlibrary{arrows.meta} %画箭头
\usetikzlibrary{commutative-diagrams} %交换图
\usetikzlibrary{positioning,automata}
\usepackage[all]{xy}
\usepackage{booktabs} %三线表
\usepackage{multirow} %合并单元格
\usepackage{makecell} %表格内强制换行
\usepackage{array}
\usepackage{booktabs} %控制表格行高
\usepackage{colortbl}  %彩色表格需要加载的宏包
\usepackage{amsmath}
\usepackage{unicode-math}
\usepackage{fontspec} %字体
\newcommand*{\cred}{\textcolor{red}}

\setmainfont{Times New Roman}
\setmathfont{XITS Math}

\setlength{\parindent}{0pt} %无缩进
\usepackage{setspace}
\usepackage{fancyhdr} %设置页码
\usepackage{lastpage} %获取总页数
\pagestyle{fancy} %fancyhdr宏包新增的页面风格
\renewcommand\headrulewidth{0pt} %取消页眉中的装饰分割线
\fancyhf{}
\cfoot{\footnotesize 第 \thepage\ 页(共 \pageref{LastPage}页)}%当前页 of 总页数

\begin{document}
\begin{spacing}{1.625}

\centerline{{\huge 河北农业大学课程考试试卷}}

\centerline{\large 2023-2024学年第\underline{\ 1\ }学期}

考试科目：\underline{物理化学与胶体化学} \quad 姓名：\underline{ \qquad \qquad} \quad 学号：\underline{ \qquad \qquad \qquad}

学院：\underline{理工学院} \quad 专业班级：\underline{\qquad \qquad \qquad \qquad} \quad 卷别：\cred{B} 考核方式：\cred{开卷}

\bigskip


一、思考题(共40分)

1. (5分) 25℃，101.325KPa下，液态氮的标准摩尔生成焓为零吗？原因？

\bigskip
2. (5分) 焓是系统以热的方式交换的能量，这一说法是否正确，为什么？

\bigskip
3. (5分) 在标准压力下，将室温下的水蒸发为同温同压的气，如何设计可逆过程？

\bigskip
4. (5分) 体系中含有\ch{H2O、H2SO4、H2SO4. 4 H2O、H2SO4. 2 H2O、H2SO4. H2O}，其组分数C是多少？说明原因。

\bigskip
5. (5分) 可逆电池与不可逆电池有哪些区别？

\bigskip
6. (5分) 有人为反应\ch{Na + H2O -> NaOH (aq) + 1/2 H2}设计了如下电池：

\ch{Na|NaOH(aq)|H2|Pt}。此电池设计错在哪里？

\bigskip
7. (5分) 简述实验测定反应\ch{A + B -> R}的反应级数的方法。

\bigskip
8. (5分) 把一个大液滴和一个小液滴密封在一个玻璃罩内，隔相当长时间后，估计会出现什么现象？并解释出现该现象的原因。

\newpage
二、计算题(共60分)

1.(20分) 0.1 mol单原子理想气体，始态为400K, 101.325 kPa, 经下列两个途径到达相同的终态：

(1) 等温可逆膨胀至\SI{10}{dm^{-3}}，再等容升温至610K。

(2) 绝热可逆膨胀至\SI{6.56}{dm^{-3}}，再等压升温至610K。

请分别计算两个途径的$Q$，$W$，$\Delta U$和$\Delta H$。若只知始态和终态，能否求出两途径的$\Delta U$和$\Delta H$？

\vspace{4em}
2. (15分) 
蔗糖\ch{C12H22O11 (s)}在人体中最终代谢产物为\ch{CO2 (g)}和\ch{H2O (l)}，求1 mol蔗糖在人体中(37℃)完全氧化时的反应熵变。已知蔗糖在298K时的定压摩尔热容$C_{p,\mathrm{m}}=425.51$ \si{J.mol^{-1}.K^{-1}}。

\vspace{4em}
3. (10分)
在25℃时，将某电导池装入\SI{0.1}{mol/L}的KCl溶液，测得电阻为23.78 $\Omega$。换成\SI{0.002414}{mol/L}的HAc溶液，测得电阻为3942 $\Omega$。求HAc溶液的电离度。

\vspace{4em}
4. (15分)
溴乙烷分解反应的活化能$E_\mathrm{a}=229.3$ \si{kJ.mol^{-1}}，650 K时，速率常数$k=$\SI{2.14e-4}{s^{-1}}，要使该反应在10 min内完成90\%，反应温度应控制在多少？


\end{spacing}
\end{document}
